\begin{titlepage}
  \centering
  \vspace*{8mm}
  {\bfseries\fontsize{26}{31}\selectfont BrainChip AKD1000\par}
  \vspace{5mm}
  {\fontsize{16}{21}\selectfont
  Manual de puesta en marcha, ejecución de modelos y diagnóstico\par}
  \vspace{8mm}
  \includegraphics[width=0.60\textwidth]{hardware-powered-isolated.png}
  \vfill
  {\large Álvaro Poveda Ruiz\par}
  \vspace{2mm}
  {\small Manual de usuario\par
  Universidad Politécnica de Madrid, 2026\par}
  \vspace{8mm}
  {\small Versión \versionmanual{} · Revisión: \fechacorte\par}
\end{titlepage}

\chapter*{Resumen}
\addcontentsline{toc}{chapter}{Resumen}

Este manual describe la puesta en marcha de una tarjeta BrainChip AKD1000 conectada por PCIe.
El procedimiento incluye la preparación del equipo, la instalación del controlador y del
entorno de Python, la carga de un modelo y la comprobación de su ejecución. También explica
cómo medir la latencia y localizar los fallos más habituales.

La primera prueba utiliza un modelo sencillo cuya salida se conoce de antemano. Se ejecuta
primero en software y después en la tarjeta. Así se puede comprobar por separado el
funcionamiento del entorno de Akida, la compatibilidad del modelo y el acceso al hardware
antes de utilizar una aplicación propia.

La ejecución completa de esta versión del procedimiento en una placa física sigue pendiente.

\section*{Alcance del manual}

El manual cubre las operaciones necesarias para instalar el entorno, detectar la tarjeta,
cargar y ejecutar modelos, consultar las estadísticas y medir el rendimiento.
También incluye comprobaciones para localizar errores de conexión, controlador, datos o modelo.

Las instrucciones de montaje deben aplicarse a la revisión concreta de cada componente.
Antes de conectar la alimentación, hay que comprobar la tensión, la polaridad y la orientación
del cable en la documentación del fabricante. El manual no cubre la reparación electrónica
de la tarjeta ni su diseño interno.

\chapter*{Cómo usar este documento}
\addcontentsline{toc}{chapter}{Cómo usar este documento}

Descarga el PDF junto con \code{akd1000-bringup.zip} y extrae el ZIP. Abre una terminal
en la carpeta \projectroot, donde están \code{pyproject.toml}, \code{examples/} y
\code{models/}. Todos los comandos del proyecto parten de esa carpeta, salvo que se indique otra.
No ejecutes los ejemplos desde el visor del ZIP. Conserva el PDF y el ZIP con la misma versión.

Ejecuta solo el bloque que corresponda a tu equipo.
Los bloques que usan \code{sudo}, \code{apt} o \code{lspci} son para Linux. El
\cref{chap:software} incluye la activación del entorno en Windows.

Si trabajas sin tarjeta, empieza por el \cref{chap:software} y completa la prueba en software.
Si dispones del hardware, lee primero el inventario y el montaje de los capítulos
\ref{chap:orientacion} y \ref{chap:hardware}; después prepara el controlador y el SDK.

\section*{Versión y material incluido}

\begin{tabularx}{\textwidth}{@{}L{0.29\textwidth}Y@{}}
\toprule
\textbf{Campo} & \textbf{Valor} \\
\midrule
Versión y fecha & \versionmanual{}; \fechacorte \\
Fuente editable & \code{manual/main.tex} y \code{manual/sections/} \\
Pruebas de esta edición & Pruebas automatizadas y simulador con Akida 2.19.1 y 2.19.3 en Windows; ejecución física pendiente \\
Modelo incluido & \code{models/akida\_v1\_smoke\_test.fbz}; cuatro entradas y dos salidas \\
\bottomrule
\end{tabularx}

Las referencias finales contienen enlaces a las instrucciones oficiales. La ruta de software
actual y la compatibilidad del controlador se comprueban por separado: que un SDK se instale
no garantiza que el kernel pueda cargar el controlador PCIe.
